\section{Conventions and the three hypotheses}\label{sec:premises}

All spaces are over \(\C\), unless a statement in the projective
appendices explicitly specifies another field of characteristic zero.
A compact K\"ahler space means a complex analytic space with a
K\"ahler form given by local smooth strictly plurisubharmonic potentials
on local embeddings. Our spaces are normal and connected, hence
irreducible. Smooth compact K\"ahler resolutions and resolutions of
meromorphic maps are obtained by projective modifications. We use the
usual resolution and K\"ahler modification theorems
\cite{HironakaResolution,VarouchasKahler}.

\subsection{Actual adjoints and positivity}

The canonical object on a normal space \(X\) is its reflexive canonical
sheaf \(\omega_X\). For an effective rational Weil divisor \(B\), the
notation \(D=K_X+B\) is an actual rational holomorphic line bundle
when, for some positive integer \(r\) clearing the coefficients of
\(B\),
\[
 L_r=\bigl(\omega_X^{[r]}\otimes\OO_X(rB)\bigr)^{**}
\]
is invertible. Its positive tensor powers define \(\OO_X(krD)\).
Here \(\omega_X^{[r]}=(\omega_X^{\otimes r})^{**}\), and divisorial
sheaves and their products are interpreted reflexively. An identity
\(D_1\sim_{\Q}D_2\) means an isomorphism of actual holomorphic line
bundles after a common positive integral multiple. Numerical equivalence
alone is never used to infer this identity.

Canonical divisor calculations can be performed with compatible local
canonical divisors on a common smooth model \(p:W\to X\). We use log
discrepancies
\begin{equation}\label{eq:discrepancy}
 a(E;X,B)=1+\operatorname{coeff}_E
       \bigl(K_W-p^*(K_X+B)\bigr).
\end{equation}
The pair is klt if these numbers are positive for every prime divisor
over \(X\), including primes on \(X\). In particular the boundary
coefficients are less than one. We do not assume at the outset that a
positive canonical power has a nonzero meromorphic section.

The adjoint \(D\) is \emph{analytically pseudo-effective} if
\(c_1(L_r)/r\) is represented by a closed positive \((1,1)\)-current
with local potentials. It is \emph{analytically nef} if this class is
in the closure of the K\"ahler cone in real Bott--Chern cohomology.
Equivalently, for a fixed K\"ahler form \(\omega\) and every
\(\varepsilon>0\), the line bundle has a smooth Hermitian metric
whose curvature, divided by \(r\), is bounded below by
\(-\varepsilon\omega\). On singular spaces the forms and potentials
are understood through local embeddings. Nonnegativity on compact
curves is a consequence of nefness, not its definition.

The Iitaka dimension \(\kappa(X,D)\) is \(-\infty\) if all positive
Cartier multiples have zero sections. Otherwise it is the maximum
dimension of the images of their complete linear systems. A rational
line bundle is \emph{semiample} if some positive Cartier multiple is
globally generated. In particular, a semiample rational line of Iitaka
dimension zero is torsion as an actual rational line bundle.
We write \(a(X)\) for algebraic dimension, and
\(q(M)=h^1(M,\OO_M)\) on smooth compact K\"ahler models.

\begin{definition}[Global strong \(\Q\)-factoriality]
\label{def:strong}
A normal compact space \(X\) is globally strongly
\(\Q\)-factorial if every coherent rank-one reflexive sheaf
\(\cF\) on the whole \(X\) has an invertible positive reflexive
power \(\cF^{[m]}=(\cF^{\otimes m})^{**}\).
\end{definition}

Definition~\ref{def:strong} is a condition on global sheaves. It does
not assert the same condition on every analytic open subset. Smooth
spaces satisfy it. The constructions below retain this global
condition when it is required; local analytic factoriality is not
inserted as an intermediate assumption.

\subsection{Orbifold subadditivity}

We give the model convention in the first hypothesis because the
distinction between the invariant base and the base on an arbitrary
model is used in the proof.

For a smooth compact space \(Z\) with an effective rational SNC
boundary \(\Gamma\) having coefficients in \([0,1]\), set
\[
 m_\Gamma(E)=
 \begin{cases}
 (1-\operatorname{coeff}_E\Gamma)^{-1},&
                  \operatorname{coeff}_E\Gamma<1,\\
 \infty,&\operatorname{coeff}_E\Gamma=1
 \end{cases}
\]
for every prime divisor \(E\). Thus the multiplicity off the boundary
is one; finite multiplicities need not be integers. If
\(g:(Z,\Gamma)\to S\) is a surjective morphism with connected
fibers and smooth base, put
\begin{equation}\label{eq:orbifold-base}
 \begin{split}
 m(g,\Gamma;P)&=
   \min_{E:\,g(E)=P}\bigl(\ord_E(g^*P)m_\Gamma(E)\bigr),\\
 B(g,\Gamma)&=\sum_P\left(1-\frac1{m(g,\Gamma;P)}\right)P.
 \end{split}
\end{equation}
Only divisors dominating \(P\) enter the minimum; the convention is
\(1/\infty=0\). The sum has finite support. This is the
inf-multiplicity convention.

An elementary equivalence of smooth source/base fibrations is a
commuting bimeromorphic diagram
\[
\begin{tikzcd}
 (Z',\Gamma')\arrow[r,"p"]\arrow[d,"g'"']&
 (Z,\Gamma)\arrow[d,"g"]\\
 S'\arrow[r,"q"]&S,
\end{tikzcd}
\]
where the maps are proper holomorphic modifications,
\(p_*\Gamma'=\Gamma\), and \(p\) is an orbifold morphism: for
every prime \(D\) on \(Z\) and every prime \(E\) on \(Z'\) with
\(t=\ord_E(p^*D)>0\),
\begin{equation}\label{eq:orbifold-morphism}
 t\,m_{\Gamma'}(E)\ge m_\Gamma(D).
\end{equation}
Infinite multiplicities are compared in the usual order, with
\(\infty\ge\infty\). All boundaries in these diagrams are rational
SNC boundaries in \([0,1]\). Equivalence is generated by these
diagrams with arrows allowed in either direction, not by arbitrary
boundary changes on a fixed space. Define
\begin{equation}\label{eq:orbifold-invariant}
 \kappa(g\mid\Gamma)=
 \inf_{g'\sim g}\kappa\bigl(S',K_{S'}+B(g',\Gamma')\bigr).
\end{equation}
For an initially normal possibly singular base, first resolve that
base and the main fiber product. On the resulting smooth source use
the strict transform of the original boundary plus the full reduced
exceptional divisor, and resolve the total support to SNC. These
modifications are isomorphisms over very general base points. The
invariant is independent of these choices.

A smooth-base fibration is \emph{neat} if there is a proper
bimeromorphic orbifold morphism of its source pair to a smooth pair,
with the boundary pushing forward, which contracts every source prime
mapping to codimension at least two in the given base. The hypothesis
below includes existence of suitable neat models and the formula
\begin{equation}\label{eq:neat-formula}
 \kappa(g\mid\Gamma)=\kappa\bigl(S,K_S+B(g,\Gamma)\bigr)
 \qquad\text{on a neat model}.
\end{equation}

\begin{assumption}[Full orbifold Iitaka subadditivity]
\label{hyp:campana}
Let \(X\) be a smooth compact connected manifold in Fujiki class
\(\mathcal C\), let \(\Delta\) be an effective rational SNC
boundary with coefficients in \([0,1]\), and let \(f:X\to Y\) be
a surjective holomorphic map with connected fibers onto a normal
compact irreducible complex space. In arbitrary dimensions,
\begin{equation}\label{eq:campana}
 \kappa(X,K_X+\Delta)\ge
 \kappa(F,K_F+\Delta|_F)+\kappa(f\mid\Delta),
\end{equation}
where \(F\) is a very general smooth fiber with SNC boundary
restriction, and the invariant base is defined by
\eqref{eq:orbifold-base}--\eqref{eq:neat-formula}.
\end{assumption}

Here very general means outside a countable union of proper closed
analytic subsets, as well as the critical values and unsuitable
boundary-stratum loci. The zero boundary is allowed, a point has
Kodaira dimension zero, and a sum with \(-\infty\) is
\(-\infty\). Neither source nor base need be projective. There is
no abundance, positivity or good-model premise in
Assumption~\ref{hyp:campana}. The displayed statement is the orbifold
subadditivity theorem of \cite[Theorem~1.1]{GrantCampana}. For later use,
on a smooth base the invariant
is at least \(\kappa(Y,K_Y)\), by effectivity of the orbifold
boundaries and smooth birational invariance.

\subsection{The pseudo-effective fourfold program}

The closed analytic cone \(\NA(X)\) consists of positive closed
bidimension-\((1,1)\) currents, modulo their pairings with all real
Bott--Chern \((1,1)\)-classes. We use its negative extremal rays in
the following hypothesis.

\begin{assumption}[Pseudo-effective fourfold MMP]
\label{hyp:mmp}
Assume Assumption~\ref{hyp:campana}. Start from an ordinary klt pair
\((X,B)\), where \(X\) is a normal irreducible globally strongly
\(\Q\)-factorial compact K\"ahler fourfold, \(B\) is effective
and rational, and the actual adjoint \(D=K_X+B\) is
\(\Q\)-Cartier and analytically pseudo-effective. No initial
modification is inserted. The following program exists and
terminates.

At every non-nef stage \((X_i,B_i)\) there is a nonzero
\(D_i\)-negative extremal ray of \(\NA(X_i)\), and every such
ray \(R\) may be chosen. It has a nef supporting class \(\alpha\)
with
\[
 \NA(X_i)\cap\alpha^\perp=R,
 \qquad \alpha-c_1(D_i)\text{ K\"ahler}.
\]
For every chosen ray there is a projective surjective bimeromorphic
contraction with connected fibers
\(f_i:X_i\to Z_i\), where \(Z_i\) is normal compact K\"ahler,
\(\rho_{\BC}(X_i/Z_i)=1\), and \(-D_i\) is relatively ample.
For some K\"ahler class \(\omega_i\) on \(Z_i\),
\[
 \NA(X_i)\cap(f_i^*\omega_i)^\perp=R.
\]
A divisorial contraction gives the next pair by pushforward. For a
small contraction the canonical log-canonical-positive flip is the
relative analytic Proj of
\[
 \bigoplus_{m\ge0}(f_i)_*\OO_{X_i}(mrD_i)
\]
for a sufficiently divisible positive Cartier index \(r\). This
algebra is locally finitely generated. Its Proj is normal and its
map to \(Z_i\) is projective, small and has connected fibers. The
boundary is strictly transformed; the new actual adjoint is
relatively ample, and its divisible Cartier multiple is the
tautological line bundle. Passing to a Veronese does not change the
model.

Each new pair is klt, compact K\"ahler and globally strongly
\(\Q\)-factorial, with \(\Q\)-Cartier analytically
pseudo-effective actual adjoint. Every nonterminal finite prefix
can be extended. Every program so obtained has finitely many steps,
irrespective of the choices of negative rays. It ends at an
analytically nef pair. The composite map extracts no prime and
satisfies the minimal-model discrepancy inequalities, weakly for
every prime over the models and strictly for every original prime
contracted.
\end{assumption}

Assumption~\ref{hyp:mmp} is supplied by
\cite[Theorem~6.20; Proposition~3.4 and its proof]{GrantMMP}
in the stated pseudo-effective subcase, without an orbifold Iitaka
hypothesis. The assumption gives neither a first plurisection
nor semiampleness. It includes no non-pseudo-effective program and
no dlt or generalized-pair extension. Fiber-type contractions are
not stopping outcomes in this program. The Cartier index may change
at every stage; no positive-side relative Bott--Chern dimension or
identification of absolute Bott--Chern spaces across a flip is
assumed. Trivial flops, inserted blowups and arbitrary birational
walks are not steps of the program. These distinctions matter when
auxiliary boundaries are used below.

\subsection{Abundance after nonvanishing}

\begin{assumption}[Effective fourfold abundance]
\label{hyp:effective-abundance}
Let \(X\) be a normal connected compact K\"ahler fourfold and
\(\Delta\) an effective rational boundary such that
\((X,\Delta)\) is klt and the actual adjoint
\(J=K_X+\Delta\) is \(\Q\)-Cartier. If \(J\) is analytically
nef and \(\kappa(X,J)\ge0\), then a positive Cartier multiple of
\(J\) is globally generated.
\end{assumption}

Assumption~\ref{hyp:effective-abundance} is the abundance-after-nonvanishing
theorem of \cite[Theorem~1.1]{GrantEffectiveAbundance}. It needs no strong
factoriality hypothesis and includes the actual torsion conclusion
when \(\kappa=0\). Its nonvanishing premise is explicit. No
auxiliary statement from its proof is granted separately.

The proof uses all three hypotheses. Apart from them, established
literature is invoked at its stated scope. The separate projective
abundance theorem used in the Moishezon branch is proved in
Appendices~\ref{proj:sec:introduction}--\ref{proj:sec:completion};
Lemma~\ref{red:log-iitaka} establishes its sole subadditivity premise
from Assumption~\ref{hyp:campana}.
